\documentclass[12pt]{article}
\usepackage[a4paper,margin={1.0cm, 2.0cm}]{geometry}
\usepackage{xargs}
\usepackage{amsmath,amssymb}
\usepackage{hyperref}
\usepackage{physics}
\usepackage{graphicx}

% signature
\newcommand{\sgn}{\text{sgn}}
% index sequence
\newcommand{\seq}[1]{\langle #1\rangle}
% ascending sequence
\newcommand{\asc}[1]{\upharpoonleft #1 \upharpoonright}
% hermitian conjugate
\newcommand{\hc}{^\dagger}
% inverse
\newcommand{\inv}{^{-1}}
\newcommand{\Sym}{\text{Sym}}
% normal ordering
\newcommand{\normord}[1]{:\mathrel{\mspace{2mu}#1\mspace{2mu}}:}
% Expectation value
\newcommand{\exval}[1]{\mathbb{E}\left[ #1 \right]}
% low-excitation projection
\newcommand{\lep}{_{\text{LE}}}
% diagonalisation of vector
\newcommand{\diag}[1]{\text{diag}\left(#1\right)}

\newcommand*\chem[1]{\ensuremath{\mathrm{#1}}}

\newtheorem{theorem}{Theorem}[section]
\newtheorem{lemma}[theorem]{Lemma}

\begin{document}

	\title{Potential energy surface extrapolation using adiabatic deformation of $SU(M)$ coherent states}
	\author{Michal Horanský}	
	\maketitle
	
	\abstract{Diagonalisation of a coherent-state sample around the HF reference state is proposed as a post-HF method to estimate the molecular-electronic ground state. The $SU(M)$ fermionic Thouless states are described with regard to the relation between their parameter matrix and their overlap with single-excitation Slater determinants. For a RHF basis, a method which guides the sampling process based on the diagonalisation on zero- and single-excitation Slater determinant subspace is described and demonstrated on spin-zero molecules \chem{Li_2} and \chem{N_2}. Generalisation of the sampling guidance process for an UHF basis is described and demonstrated on spin-one molecules \chem{BeH} and \chem{NO}.}
	
	%\newpage
	
	\tableofcontents
	
	\section{Background}
	
	\subsection{Fermionic $SU(M)$ Thouless states in electronic structure}
	
	Consider a Hilbert space described by $M$ modes occupied in total by $S$ fermions, where $S < M$. The occupancy basis of this Hilbert space consists of the Slater determinants
	
	$$\ket{n_1, n_2 \dots n_M} \qq{where} \sum_{i=1}^M n_i = S \qq{and} n_i \in \{0,1\}$$
	
	By defining the reference state as
	
	\begin{equation}
	\ket{\phi_0} = \ket{1, 1 \dots 1, 0, 0 \dots 0}
	\end{equation}
	
	we may label every occupancy basis state as $\ket{\phi_0 + a - b}$, where $a, b$ are equal-sized sets of indices with the constraints
	\begin{eqnarray*}
	a_i \in a&\qq{:}& 1 \leq a_i \leq M - S\\
	b_i \in b&\qq{:}& 1 \leq b_i \leq S
	\end{eqnarray*}
	and the state $\ket{\phi_0 + a - b}$ describes the Slater determinant where modes occupied in $\phi_0$ are occupied iff their index is not in $b$, and modes unoccupied in $\phi_0$ are occupied iff their index is in $a$ (with the first unoccupied mode, i.e. the $S+1$-st mode, having index $1$).
	
	It will be helpful to define the integer sets
	\begin{eqnarray*}
	\pi_1 &=& \{1, 2 \dots S\}\\
	\pi_0 &=& \{1, 2 \dots M - S\}
	\end{eqnarray*}
	which enumerate the occupied and unoccupied modes in $\ket{\phi_0}$, respectively.
	
	Then, we can construct the $SU(M)$ Thouless states restricted to the Hilbert space spanned by the aforementioned basis from the reference state like so:
	
	\begin{equation}
	\ket{Z} = N(Z, Z^*) \exp(\sum_{i\in\pi_0}\sum_{j\in\pi_1}\hat{f}_{i+S}\hc \hat{f}_j) \ket{\phi_0} 
	\end{equation}
	
	where $N$ is the normalisation factor and $\hat{f}\hc, \hat{f}$ are the fermionic creation and annihilation operators, respectively. Note that $Z$ is therefore a $(M-S) \cross S$ matrix of complex number.
	
	These states form a non-orthogonal, overcomplete basis of the full Hilbert space.
	
	It is then a known result that the decomposition of this Thouless state into the occupancy basis is
	
	\begin{equation} \label{eq:decomposition}
	\braket{\phi_0 + a - b}{Z} = N(Z, Z^*) (-1)^{\frac{1}{2}|a|(|a| + 1)} \det(Z_{a,b})
	\end{equation}
	where $|a|=|b|$ is the number of excitations from the reference state and $Z_{a,b}$ is the submatrix of $Z$ composed of rows and columns labelled by indices in $a$ and $b$, respectively.
	
	In other words, the minors of $Z$ describe the decomposition coefficients of specific occupancy basis states in $\ket{Z}$, with all minors of order $r$ describing all occupancy basis states obtained by promoting exactly $r$ fermions from modes occupied in $\ket{\phi_0}$ into modes unoccupied in $\ket{\phi_0}$.
	
	It is also known that the overlap between two states is
	
	\begin{equation}
	\braket{Z_a}{Z_b} = N(Z_a, Z_a^*)N(Z_b, Z_b^*) \det(I + Z_a\hc Z_b)
	\end{equation}
	
	which also determines the normalisation factor
	\begin{equation}
	N(Z, Z^*) = \frac{1}{\sqrt{\det(I + Z\hc Z)}}
	\end{equation}
	
	\subsection{Electronic structure described by pairs of Thouless states}
	
	Under the Born-Oppenheimer approximation, we can describe the wavefunction of the electrons of a particular molecule by a vector in a Hilbert space spanned by Slater determinants constructed on molecular spin-orbitals:
	
	\begin{equation}
	\ket{\Psi} = \ket{n(\phi_1^\alpha), n(\phi_1^\beta), n(\phi_2^\alpha), n(\phi_2^\beta) \dots}
	\end{equation}
	where, for practicality, we terminate the sequence of molecular orbitals at some highest orbital $\phi_M$, where we hope that the behaviour of the molecule in the energy range of interest as described on this finite basis closely matches its true behaviour.
	
	As the labelling suggests, it is very natural to decompose this Hilbert space into a direct product of (spatial) molecular orbitals and the two spin states $\ket{\alpha}, \ket{\beta}$, which, assuming the Hamiltonian has full rotatinal symmetry, can be selected arbitrarily as any orthonormal basis of the spin-$\frac{1}{2}$ system (typically, these are associated with spin-up and spin-down states, i.e. eigenstates of $\hat{S}_z$ for the particular selection of the $z$-axis). Then, the occupancy basis vectors may be writ out like so:
	
	\begin{equation}
	\ket{\Psi} = \ket{n_1, n_2 \dots n_M}_\alpha \otimes \ket{\bar{n}_1, \bar{n}_2 \dots \bar{n}_{\bar{M}}}_\beta
	\end{equation}
	
	where $n_i, \bar{n}_j$ are shorthands for the occupancy numbers of modes $\phi_i^\alpha, \phi_j^\beta$, respectively. Note that we do not implicitly assume that $\phi_i^\alpha = \phi_i^\beta$, as is seen by allowing even the total number of the MOs in each spin-subspace to be different. The case where this equivalence is assumed is referred to as a restricted Hartree-Fock (RHF) basis. Otherwise, we are working under the unrestricted Hartree-Fock (UHF) basis.
	
	Note that in the non-relativistic case the Hamiltonian preserves the total number of electrons in each spin-subspace:
	\begin{equation}
	[\hat{H}, \hat{N}_\alpha] = [\hat{H}, \hat{N}_\beta] = 0 \qq{where} \hat{N}_\alpha \ket{\Psi} = \sum_{i=1}^M n_i = N_\alpha \qq{and} \hat{N}_\beta \ket{\Psi} = \sum_{i=1}^{\bar{M}} \bar{n}_i = N_\beta
	\end{equation}
	Equivalently, $\hat{H}$ commutes with the total number operator $\hat{N}=\hat{N}_\alpha + \hat{N}_\beta$ and the total spin $z$-projection operator $\hat{S}_z = \hat{N}_\alpha - \hat{N}_\beta$. As a result, $N_\alpha, N_\beta, S_z$ are good quantum numbers.
	
	Consequently, the two spin-subspaces of the full Hilbert space are spanned exactly by the kind of basis which admits the construction of fermionic $SU(M)$ Thouless states. Hence, we can span the Hilbert space by states of the form
	
	\begin{equation}
	\ket{Z^\alpha, Z^\beta} = \ket{Z^\alpha}_\alpha \otimes \ket{Z^\beta}_\beta
	\end{equation}
	where $\ket{Z^\alpha}, \ket{Z^\beta}$ are Thouless states with mode and total occupancy numbers being $M, N_\alpha$ and $\bar{M}, N_\beta$, respectively.
	
	\textit{Note.} It will be helpful to denote
	\begin{equation}
	\ket{(B \rightarrow A), (B' \rightarrow A')}
	\end{equation}
	to be the occupancy basis state obtained by promoting electrons from modes $B$ to modes $A$ in the spin-alpha subspace and electrons from modes $B'$ to modes $A'$ in the spin-beta subspace. For a single promotion, the indices of particular modes replace the index sets.
	
	\section{Decomposing the FCI vector into the coherent state basis} \label{sec:FCI to CS}
	
	This is a tricky decomposition. Since it is a discrete basis sample, it will not span the entire Hilbert space, and we must be careful to approximate the FCI vector well. We want the basis to be small, which means conditioning the stochastic sampling to optimise the process.
	
	Ideas for sampling:
	\begin{itemize}
	\item \textbf{Gradient descent/annealing}. We take some reference point, descend/anneal to find a CS with large overlap, then subtract it from the FCI and repeat on the reduced vector to whittle it down.
	
	 \textbf{Problems}: annealing must be used so we don't leave too many significant active subspaces; temperature is a free parameter; multivariate FCI states must be labelled by multiple reference states from which we initiate the annealing; conditioning by orthogonoality is difficult.
	 \item \textbf{Sampling from wavefunction with rejection}. We choose some PDF $F(Z)$ which should cover all parts of the Hilbert space where FCI is significant. We sample a state from $F$, then reject it with probability $F(Z)|\braket{Z}{\rm FCI}|^2$. For conditioning, we take an array of candidates; for the unrejected candidates, we choose the one which has the smallest projection norm into the existing basis.
	 
	 \textbf{Problems}: Still struggling with multivariate FCIs and $F$ must be chosen wide enough to not forget about any active spaces; this in turn slows down the sampling process by raising the rate of rejection. However, it does converge in finite time with absolute certainty.
	 
	 \textbf{\textit{This is our strongest candidate.}}
	\end{itemize}
	
	
	
	
	
	
	
	\section{Adiabatic evolution of a superposition of coherent states}
	
	The time-evolution of coherent states is a well-studied subject. Gilmore-Perelomov coherent states are useful constructs precisely because they remain approximately coherent under evolution with a Hamiltonian quadratic in the elements of the dynamical group's Lie algebra; in the case of fermionic $SU(M)$ states, these elements are $\hat{f}_i\hc \hat{f}_j$. The parameter tensor obeys the following equation of motion:
	
	\begin{equation}\label{eq: cs evolution}
	\dot{z}_i = -\frac{i}{\hbar} \xi^T_{ij}(z, z^*) \pdv{H(z, z^*)}{z_j}\qq{where} \xi^{-1}(z_a^*, z_b) = \pdv[2]{\ln{\unnormbraket{z_a}{z_b}}}{z_b}{z^*_a}\qq{and}H(z, z^*) = \frac{\unnormmel{z}{\hat{H}}{z}}{\unnormbraket{z}{z}}
	\end{equation}
	Here, $\xi$ is the inverse of the quotient-space metric.
	
	When inquiring about the time evolution about a superposition of multiple coherent states, with the goal of capturing the dynamics beyond the approximation which treats the CSs as \textit{actually} coherent, two approaches are immediately to be considered: firstly, one may treat the component coherent states separately, propagating them according to Eq.~\eqref{eq: cs evolution}, and then use variational principle to determine the dynamics of the decomposition coefficients; secondly, one may use the full variational principle to consider the cross-dependence between all dynamical components, including the parameter tensor components $z_{a,i}$ and the decomposition coefficients $c_a$.
	
	The subject of this section is to describe the analogous, but distinct behaviour of coherent states and their superpositions under the adiabatic evolution of an energy eigenstate. Analogous to the problem of time evolution, two approaches are immediately to be considered. The first propagates each CS separately and then considers the decomposition coefficients, whilst the other employs the full variational principle, coupling all dynamical variables.
	
	The desired output for each approach is the \textit{transfer function} describing the evolution $z\to z + \delta z, c \to c + \delta c$ under the deformation $\hat{H}(q = 0) \to \hat{H}(q = \delta q)$, where $\ket{\Psi (q = 0)} = \sum_a^N c_a \unnormket{z_a}$ is an energy eigenstate of $\hat{H}(q = 0)$.
	
	\subsection{Properties of the tangent space of the coherent-state manifold}
	x
	The complex parametrisation $z \to \unnormket{z}$ is in some sense easier to work with than its projection onto the norm-1 Hilbert subspace $\ket{z}$, since it has no explicit dependence on $z^*$ (lacking the normalisation factor $N(z, z^*)$). Its tangent space permits a natural choice of basis given by $\pdv{z_i}\unnormket{z}$, which can be trivially shown to be
	\begin{equation}
	\pdv{z_i}\unnormket{z} = \hat{T}_i \unnormket{z}
	\end{equation}
	from the definition of Thouless parametrisation, \textit{under the assumption} that the transition operator commutators all vanish. This is indeed the case for fermionic $SU(N)$ coherent states.
	
	Consider now the expectation value of any stationary\footnote{Commuting with $\pdv{z_i}$ and $\pdv{z_i^*}$. Note that $\hat{T}_i, \hat{T}_j\hc$ are both stationary, even if they do not necessarily commute together.} operator $\hat{A}$:
	\begin{equation}
	\seq{\hat{A}} = A(z, z^*) = \frac{\unnormmel{z}{\hat{A}}{z}}{\unnormbraket{z}{z}}
	\end{equation}
	
	In anticipation of the analysis to come, we will here derive some properties of $\seq{A}$ with regard to differentiation against $z, z^*$. Firstly, the explicit first-order derivatives:
	\begin{align}
	\pdv{z_i}\seq{\hat{A}} &= \frac{1}{\unnormbraket{z}{z}^2}\left(\unnormmel{z}{\hat{A}\hat{T}_i}{z}\unnormbraket{z}{z} - \unnormmel{z}{\hat{A}}{z}\unnormmel{z}{\hat{T}_i}{z}\right) = \seq{\hat{A}\hat{T}_i} - \seq{\hat{A}}\seq{\hat{T}_i}\\
	\pdv{z_i^*}\seq{\hat{A}} &= \seq{\hat{T}_i\hc\hat{A}} - \seq{\hat{A}}\seq{\hat{T}_i\hc}
	\end{align}
	Now the second-order derivatives:
	\begin{align}
	\pdv[2]{}{z_i}{z_j}\seq{\hat{A}} &= \frac{1}{\unnormbraket{z}{z}^4}\left( \unnormmel{z}{\hat{A}\hat{T}_i\hat{T}_j}{z}\unnormbraket{z}{z}^3 + \unnormmel{z}{\hat{A}\hat{T}_i}{z}\unnormmel{z}{\hat{T}_j}{z}\unnormbraket{z}{z}^2 - \unnormmel{z}{\hat{A}\hat{T}_j}{z} \unnormmel{z}{\hat{T}_i}{z}\unnormbraket{z}{z}^2 \right.\nonumber\\
	 &\quad\left. - \unnormmel{z}{\hat{A}}{z}\unnormmel{z}{\hat{T}_i\hat{T}_j}{z}\unnormbraket{z}{z}^2 - 2\unnormbraket{z}{z}\unnormmel{z}{\hat{T}_j}{z}\left( \unnormmel{z}{\hat{A}\hat{T}_i}{z}\unnormbraket{z}{z} - \unnormmel{z}{\hat{A}}{z}\unnormmel{z}{\hat{T}_i}{z} \right) \right) \nonumber \\
	&= \seq{\hat{A}\hat{T}_i\hat{T}_j} - \seq{\hat{A}\hat{T}_i}\seq{\hat{T}_j} - \seq{\hat{A}\hat{T}_j}\seq{\hat{T}_i} + \seq{\hat{A}}\left(2\seq{\hat{T}_i}\seq{\hat{T}_j} - \seq{\hat{T}_i\hat{T}_j}\right)\\
	\pdv[2]{}{z_i^*}{z_j^*}\seq{\hat{A}} &= \seq{\hat{T}_j\hc\hat{T}_i\hc\hat{A}} - \seq{\hat{T}_i\hc\hat{A}}\seq{\hat{T}_j\hc} - \seq{\hat{T}_j\hc\hat{A}}\seq{\hat{T}_i\hc} + \seq{\hat{A}}\left(2\seq{\hat{T}_i\hc}\seq{\hat{T}_j\hc} - \seq{\hat{T}_j\hc\hat{T}_i\hc}\right)
	\end{align}
	Note the expected behaviour of the resulting expressions: namely, they are invariant under swapping the indices $i \leftrightarrow j$, and the two expressions are related by complex conjugation for Hermitian $\hat{A}$.
	
	The mixed second-order derivative is obtained analogously, though one must take care to preserve the order of the $\hat{T}_j\hc\hat{T}_i$ terms, since these two operators do not necessarily commute:
	\begin{equation}
	\pdv[2]{}{z_j^*}{z_i}\seq{\hat{A}} = \seq{\hat{T}_j\hc \hat{A} \hat{T}_i} - \seq{\hat{A}\hat{T}_i} \seq{\hat{T}_j\hc} - \seq{\hat{T}_j\hc \hat{A}} \seq{\hat{T}_i} + \seq{\hat{A}} \left( 2\seq{\hat{T}_i}\seq{\hat{T}_j\hc} - \seq{\hat{T}_j\hc\hat{T}_i} \right)
	\end{equation}
	
	Note the expected behaviour of this expression: it is equal to itself after simultaneously swapping the indices $i \leftrightarrow j$ and taking the complex conjugate, provided $\hat{A}$ is Hermitian.
	
	\subsection{General ansatz for adiabatic transfer of composite energy eigenstates}
	
	In this section we outline the formalised approach to our family of variational methods for adiabatic transfer. Given $\ket{\Psi(q = 0)}$ as an energy eigenstate, we wish to find $\ket{\Psi(q = \delta q)}$, the new energy eigenstate of the deformed Hamiltonian. Let us for the moment disregard the problem of degeneracy.
	
	An energy eigenstate has the unique property of remaining parallel to itself (only differing by a scalar factor) when transformed by the Hamiltonian. We shall therefore request the component of $\hat{H}_{q = \delta q}\ket{\Psi(q = \delta q)}$ which is perpendicular to $\ket{\Psi(q = \delta q)}$ to be zero. This component can be found by projecting out the parallel component:
	
	\begin{equation}
	\ket{\Psi_\perp(q = \delta q)} = \left(\hat{1} - \ket{\Psi(q = \delta q)}\bra{\Psi(q = \delta q)}\right) \hat{H}_{q = \delta q} \ket{\Psi(q = \delta q)}
	\end{equation}
	
	Although we require this to be a zero vector, that is not in itself a directly useful criterion: firstly, because the number of equations would be equal to the number of basis vectors of the Hilbert space, and secondly, because a perfect transfer is in general not possible for a finite superposition $\ket{\Psi(q = 0)}$ (which is also, in practice, only an approximate eigenstate of the Hamiltonian). For these reasons, we shall instead demand the norm squared of this perpendicular component to be minimised. Let us define
	\begin{equation}
	F(q) = \braket{\Psi_\perp(q)}{\Psi_\perp(q)} = \mel{\Psi(q)}{\hat{H}_{q}^2}{\Psi(q)} - \mel{\Psi(q)}{\hat{H}_{q}}{\Psi(q)}\mel{\Psi(q)}{\hat{H}_{q}}{\Psi(q)} \geq 0 \qq{by def.}
	\end{equation}
	
	To make the method formally acknowledge its approximate nature: we are given $\ket{\Psi(q = 0)}$ for which $F(q = 0)$ is minimised with respect to the dynamical properties $x_i$ of $\ket{\Psi}$ (in our case the decomposition coefficients $c_a$ and the coherent state parameter tensors $z_a$). Being given the derivative of the Hamiltonian with respect to the deformation parameter (in practice in the form of evaluating the deformated Hamiltonian at some small, finite deformation $\delta q$), we constrict the dynamical properties to the trajectory for which $F(q)$ is still minimal. The equation of motion is therefore:
	\begin{align}
	\dv{q} \pdv{F}{x_i} = 0
	\end{align}
	Trivially, this is equivalent to the more useful form
	\begin{equation}
	\pdv{x_i} \dv{F}{q} = 0
	\end{equation}
	A useful byproduct of this ansatz is that evaluating $F$ itself gives us a measure of the method's precision: if this number grows with the growing deformation, the method is losing accuracy away from the initial point; if this number remains stable, the method remains accurate.
	
	In the second form of the equations of motion, we may choose our set of dynamical properties (depending on the method, these will be either only the decomposition coefficients or both the decomposition coefficients and the parameter tensor components), however, the derivative of $F$ only depends on the form of $\ket{\Psi(q)}$. We shall quote its form here:
	
	\begin{equation}
	\dv{F}{q} = \pdv{F}{q} + \pdv{c_b^*}{q} \pdv{F}{c_b^*} + \pdv{c_b}{q} \pdv{F}{c_b} + \pdv{z_{b,j}^*}{q} \pdv{F}{z_{b,j}^*} + \pdv{z_{b,j}}{q} \pdv{F}{z_{b,j}}
	\end{equation}
	Before expanding this, let us define the following matrix function:
	\begin{equation}
	W_{ab}(\hat{A}) = \unnormmel{z_a}{\hat{A}}{z_b}\qq{for any operator} \hat{A}
	\end{equation}
	This matrix function has the following properties:
	\begin{equation}
	\pdv{z_{m,u}^*} W_{ab}(\hat{A}) = \delta_a^m W_{ab}(\hat{T}_u\hc\hat{A}) \qquad \pdv{z_{m,u}} W_{ab}(\hat{A}) = \delta_b^m W_{ab}(\hat{A}\hat{T}_u) \qquad W_{ab}(\hat{A})^* = W_{ba}(\hat{A}\hc)
	\end{equation}
	Where it is understood not to sum over the repeated indices $a,b$ on the right-hand sides of the two expressions.
	
	Then
	\begin{align}
	F &= c_a^* W_{ab}(\hat{H}^2)c_b - \left(c_a^* W_{ab}(\hat{H})c_b\right)^2\\
	\pdv{F}{q} &= c_a^* W_{ab}\left(\pdv{\hat{H}}{q}\hat{H} + \hat{H}\pdv{\hat{H}}{q}\right)c_b - 2\left(c_a^* W_{ab}\left(\pdv{q}\hat{H}\right)c_b\right)\left(c_c^* W_{cd}(\hat{H})c_d\right)\\
	\pdv{F}{c_a^*} &= W_{ab}(\hat{H}^2)c_b - 2W_{ab}(\hat{H})c_b \left(c_c^*W_{cd}(\hat{H})c_d\right)\\
	\pdv{F}{c_a} &= c_b^*W_{ba}(\hat{H}^2) - 2c_b^*W_{ba}(\hat{H}) \left(c_c^*W_{cd}(\hat{H})c_d\right)\\
	\pdv{F}{z_{b,j}^*} &= c_b^* W_{ba}(\hat{T}_j\hc\hat{H}^2)c_a - 2\left(c_b^* W_{ba}(\hat{T}_j\hc\hat{H})c_a\right)\left(c_c^* W_{cd}\left(\hat{H}\right)c_d\right)\\
	\pdv{F}{z_{b,j}} &= c_a^* W_{ab}(\hat{H}^2\hat{T}_j)c_b - 2\left(c_a^* W_{ab}(\hat{H}\hat{T}_j)c_b\right)\left(c_c^* W_{cd}\left(\hat{H}\right)c_d\right)
	\end{align}
	where it is understood to not sum over $b$ in the last two equations, even though the index is repeated.
	
	The full derivative of $F$ can then be expanded as
	\begin{align}
	\dv{F}{q} &= c_a^* W_{ab}\left(\pdv{\hat{H}^2}{q}\right)c_b - 2\left(c_a^* W_{ab}\left(\pdv{q}\hat{H}\right)c_b\right)\left(c_c^* W_{cd}(\hat{H})c_d\right)\nonumber\\
	&\quad + \dot{c}_a^* W_{ab}(\hat{H}^2)c_b - 2\dot{c}_a^* W_{ab}(\hat{H})c_b \left(c_c^*W_{cd}(\hat{H})c_d\right) + c_a^* W_{ab}(\hat{H}^2)\dot{c}_b - 2c_a^* W_{ab}(\hat{H})\dot{c}_b \left(c_c^*W_{cd}(\hat{H})c_d\right)  \nonumber\\
	&\quad + \left[c_a^* W_{ab}(\hat{T}_j\hc\hat{H}^2)c_b - 2\left(c_a^* W_{ab}(\hat{T}_j\hc\hat{H})c_b\right)\left(c_c^* W_{cd}\left(\hat{H}\right)c_d\right) \right] \dot{z}_{a,j}^* \nonumber\\
	&\quad + \left[c_a^* W_{ab}(\hat{H}^2\hat{T}_j)c_b - 2\left(c_a^* W_{ab}(\hat{H}\hat{T}_j)c_b\right)\left(c_c^* W_{cd}\left(\hat{H}\right)c_d\right) \right] \dot{z}_{b,j} \\
	&= c_a^* c_b \left[  W_{ab}\left(\pdv{\hat{H}^2}{q}\right) + W_{ab}(\hat{T}_j\hc \hat{H}^2) \dot{z}_{a,j}^* + W_{ab}(\hat{H}^2 \hat{T}_j) \dot{z}_{b,j}\right] \nonumber\\
	&\quad -2 c_a^* c_b c_c^* c_d W_{cd}(\hat{H}) \left[ W_{ab}\left(\pdv{q}\hat{H}\right) + \left( W_{ab}(\hat{T}_j\hc \hat{H}) \dot{z}_{a,j}^* + W_{ab}(\hat{H} \hat{T}_j) \dot{z}_{b,j} \right) \right] \nonumber\\
	&\quad + W_{ab}(\hat{H}^2) \left(\dot{c}_a^* c_b + c_a^* \dot{c}_b\right) - 2 W_{ab}(\hat{H}) W_{cd}(\hat{H}) \left( \dot{c}_a^* c_b c_c^* c_d + c_a^* \dot{c}_b c_c^* c_d \right)
	\end{align}
	and its partial derivatives with respect to the possible dynamical properties:
	\begin{align}
	\pdv{c_m^*}\dv{F}{q} &= \delta^m_a c_b \left[  W_{ab}\left(\pdv{\hat{H}^2}{q}\right) + W_{ab}(\hat{T}_j\hc \hat{H}^2) \dot{z}_{a,j}^* + W_{ab}(\hat{H}^2 \hat{T}_j) \dot{z}_{b,j}\right] \nonumber\\
	&\quad -2 \left( \delta^m_a c_b c_c^* c_d + \delta^m_c c_a^* c_b c_d \right) W_{cd}(\hat{H}) \left[ W_{ab}\left(\pdv{q}\hat{H}\right) + \left( W_{ab}(\hat{T}_j\hc \hat{H}) \dot{z}_{a,j}^* + W_{ab}(\hat{H} \hat{T}_j) \dot{z}_{b,j} \right) \right] \nonumber\\
	&\quad + \delta^m_a W_{ab}(\hat{H}^2) \dot{c}_b - 2 W_{ab}(\hat{H}) W_{cd}(\hat{H}) \left(\delta^m_c \dot{c}_a^* c_b c_d + \delta^m_a \dot{c}_b c_c^* c_d + \delta^m_c c_a^* \dot{c}_b c_d \right)\\
	\pdv{c_m}\dv{F}{q} &= \delta^m_b c_a^* \left[  W_{ab}\left(\pdv{\hat{H}^2}{q}\right) + W_{ab}(\hat{T}_j\hc \hat{H}^2) \dot{z}_{a,j}^* + W_{ab}(\hat{H}^2 \hat{T}_j) \dot{z}_{b,j}\right] \nonumber\\
	&\quad -2 \left( \delta^m_b c_a^* c_c^* c_d + \delta^m_d c_a^* c_b c_c^* \right) W_{cd}(\hat{H}) \left[ W_{ab}\left(\pdv{q}\hat{H}\right) + \left( W_{ab}(\hat{T}_j\hc \hat{H}) \dot{z}_{a,j}^* + W_{ab}(\hat{H} \hat{T}_j) \dot{z}_{b,j} \right) \right] \nonumber\\
	&\quad + \delta^m_b W_{ab}(\hat{H}^2) \dot{c}_a^* - 2 W_{ab}(\hat{H}) W_{cd}(\hat{H}) \left( \delta^m_b \dot{c}_a^* c_c^* c_d + \delta^m_d \dot{c}_a^* c_b c_c^* + \delta^m_d c_a^* \dot{c}_b c_c^* \right)\\
	\pdv{z_{m,u}^*}\dv{F}{q} &= \delta_a^m c_a^* c_b \left[  W_{ab}\left(\hat{T}_u\hc\pdv{\hat{H}^2}{q}\right) + W_{ab}(\hat{T}_u\hc\hat{T}_j\hc \hat{H}^2) \dot{z}_{a,j}^* + W_{ab}(\hat{T}_u\hc\hat{H}^2 \hat{T}_j) \dot{z}_{b,j}\right] \nonumber\\
	&\quad -2 c_a^* c_b c_c^* c_d \left[ \delta_a^m W_{ab}\left(\hat{T}_u\hc \pdv{\hat{H}}{q}\right) W_{cd}(\hat{H}) + \delta_c^m W_{ab}\left(\pdv{\hat{H}}{q}\right) W_{cd}(\hat{T}_u\hc\hat{H}) \right.\nonumber\\
	&\qquad\qquad\qquad\quad + \delta_c^m W_{cd}(\hat{T}_u\hc\hat{H}) \left( W_{ab}(\hat{T}_j\hc \hat{H}) \dot{z}_{a,j}^* + W_{ab}(\hat{H} \hat{T}_j) \dot{z}_{b,j} \right) \nonumber\\
	&\qquad\qquad\qquad\quad \left.+ \delta_a^m W_{cd}(\hat{H}) \left( W_{ab}(\hat{T}_u\hc\hat{T}_j\hc \hat{H}) \dot{z}_{a,j}^* + W_{ab}(\hat{T}_u\hc\hat{H} \hat{T}_j) \dot{z}_{b,j} \right) \right] \nonumber\\
	&\quad + \delta_a^m W_{ab}(\hat{T}_u\hc\hat{H}^2) \left(\dot{c}_a^* c_b + c_a^* \dot{c}_b\right) \nonumber\\
	&\quad - 2 \left[ \delta_a^m W_{ab}(\hat{T}_u\hc\hat{H}) W_{cd}(\hat{H}) + \delta_c^m W_{ab}(\hat{H}) W_{cd}(\hat{T}_u\hc\hat{H}) \right] \left( \dot{c}_a^* c_b c_c^* c_d + c_a^* \dot{c}_b c_c^* c_d \right)\\
	\pdv{z_{m,u}}\dv{F}{q} &= \delta_b^m c_a^* c_b \left[  W_{ab}\left(\pdv{\hat{H}^2}{q}\hat{T}_u\right) + W_{ab}(\hat{T}_j\hc \hat{H}^2\hat{T}_u) \dot{z}_{a,j}^* + W_{ab}(\hat{H}^2 \hat{T}_j\hat{T}_u) \dot{z}_{b,j}\right] \nonumber\\
	&\quad -2 c_a^* c_b c_c^* c_d \left[ \delta_b^m W_{ab}\left(\pdv{\hat{H}}{q}\hat{T}_u\right) W_{cd}(\hat{H}) + \delta_d^m W_{ab}\left(\pdv{\hat{H}}{q}\right) W_{cd}(\hat{H}\hat{T}_u) \right.\nonumber\\
	&\qquad\qquad\qquad\quad + \delta_d^m W_{cd}(\hat{H}\hat{T}_u) \left( W_{ab}(\hat{T}_j\hc \hat{H}) \dot{z}_{a,j}^* + W_{ab}(\hat{H} \hat{T}_j) \dot{z}_{b,j} \right) \nonumber\\
	&\qquad\qquad\qquad\quad \left.+ \delta_b^m W_{cd}(\hat{H}) \left( W_{ab}(\hat{T}_j\hc \hat{H}\hat{T}_u) \dot{z}_{a,j}^* + W_{ab}(\hat{H} \hat{T}_j\hat{T}_u) \dot{z}_{b,j} \right) \right] \nonumber\\
	&\quad + \delta_b^m W_{ab}(\hat{H}^2\hat{T}_u) \left(\dot{c}_a^* c_b + c_a^* \dot{c}_b\right) \nonumber\\
	&\quad - 2 \left[ \delta_b^m W_{ab}(\hat{H}\hat{T}_u) W_{cd}(\hat{H}) + \delta_d^m W_{ab}(\hat{H}) W_{cd}(\hat{H}\hat{T}_u) \right] \left( \dot{c}_a^* c_b c_c^* c_d + c_a^* \dot{c}_b c_c^* c_d \right)
	\end{align}
	
	\subsection{$U(1)$ gauge symmetry for the decomposition coefficients}
	
	For any ansatz which treats $c$ as a dynamic variable (which are all of them), we encounter a genuine singularity in the equations of motion. This is necessary, because the condition $F = 0$ does not uniquely specify the transfer function, as $F$ possesses the following gauge symmetry:
	\begin{equation}
	c \to e^{i\theta} c \qq{for} \theta \in \mathbb{R}
	\end{equation}
	Since the space of all $F$-equivalent vectors $c$ is 1-dimensional, the difference between the dimension and rank of matrix $\Omega$ in the equation of motion $\Omega \dot{c} = R$ is 1, and the space of solutions to this equation will lie on a line parametrised as $v_0 + tv_k$, where $tv_k$ is the kernel of $\Omega$.
	
	We can treat this singularity by replacing one of the rows of $\Omega$ by a gauge condition. It is important to identify a row which is linearly dependent on the other rows, since this is not necessarily true for every row in the system.
	
	In general, a gauge choice corresponds to a restriction on $t$, i.e. the component of $\dot{c}$ parallel to the kernel. The most robust and computationally safe option is to find the kernel and thus reduce $\Omega$. However, since $\Omega$ is constructed from real physical parameters (one- and two-electron exchange integrals), we argue that choosing an arbitrary, easily constructible $\dim\Omega - 1$-dimensional hyperplane to which we constrict the solution will be sufficient with near-one probability. A natural choice is to require the imaginary or real part of the inner product $c_a^* \dot{c}_a$ to vanish. Our proposed gauge choice will set its imaginary part to zero:
	
	\begin{equation}
	c_a^* \dot{c}_a - \dot{c}_a^* c_a = 0
	\end{equation}
	Should $\Omega$ be near-singular in the actual calculation, we instead set the real part to zero by changing the difference to a sum, thus rotating the hyperplane in the complex vector space.
	
	Other singularities may arise due to pathologic behaviour of the coherent-state manifold, namely when $F$ is invariant under some transformation of $z_{a,i}$ at $q = 0$. This is also a practically-zero probability case which can be ignored during actual calculation, assuming $\ket{\Psi(q = 0)}$ being non-degenerate. For the degenerate case, wait for the next article.
	
	
	\subsection{Adiabatic transfer using characteristic trajectories}
	
	In this section, we propagate each component coherent state $\unnormket{z_a}$ separately and then employ the variational principle to consider the dynamics of the decomposition coefficients $c_{a = 1\dots N}$. Since the trajectory of each coherent state is determined only by its parameter tensor $z$, a particular deformation of the Hamiltonian along a curve $\mathcal{Q}$ will map each initial value of the coherent state parameter tensor $z$ onto its characteristic trajectory $\mathcal{Z}$. For an infinetisimal deformation of the Hamiltonian, this gives us the map from the tangent bundle of the deformation parameter space\footnote{Assuming the chosen parameter space $\{q\}$ is a vector space.} onto the tangent bundle of the coherent state manifold. This map shall be referred to as the \textit{characteristic adiabatic deformation} of the coherent state manifold under the deformation parameter space.
	
	\subsubsection{Characteristic trajectories}
	
	Suppose the initial Hamiltonian $\hat{H}(q = 0) \equiv \hat{H}_0$ has an eigenstate equal to a pure coherent state $\ket{z(q = 0)}\equiv \ket{z_0}$. Then $z_0$ is stationary under time evolution (as the eigenstate remains stationary up to a phase factor), which reduces Eq.~\eqref{eq: cs evolution} to
	
	\begin{equation}
	\xi^T_{ij}(z, z^*) \pdv{H_{q = 0}(z, z^*)}{z_j}\eval_{z = z_0} = 0
	\end{equation}
	
	Because $\xi_{ij}^T(z, z^*)$ is invertible, the condition reduces to
	
	\begin{equation} \label{eq: char traj init stat}
	\pdv{H_{q = 0}(z, z^*)}{z_i}\eval_{z = z_0} = 0
	\end{equation}
	
	Under the deformation $\hat{H}(q = 0)\to \hat{H}(q = \delta q)$, we wish to find a differential $\delta z$ such that $z(q = 0) + \delta z$ satisfies the stationary condition for the deformed Hamiltonian:
	
	\begin{equation}
	\pdv{H_{q = \delta q}(z, z^*)}{z_i}\eval_{z = z_0 + \delta z} = 0
	\end{equation}
	
	Taylor expanding this condition to the first order yields
	
	\begin{equation}
	\pdv{H(z, z^*)}{z_i} + \left(\delta z_j^* \pdv{z_j^*} + \delta z_j \pdv{z_j}\right) \pdv{H(z, z^*)}{z_i} = -\pdv[2]{H(z, z^*)}{q}{z_i}\delta q
	\end{equation}
	
	where the expression is understood to be evaluated at $z = z_0, q = 0$. The first term vanishes by Eq.~\eqref{eq: char traj init stat}, leaving the equations of motion in the form
	\begin{equation}
	\pdv[2]{H(z, z^*)}{z_i}{z_j^*} \pdv{z_j^*}{q} + \pdv[2]{H(z, z^*)}{z_i}{z_j} \pdv{z_j}{q} = -\pdv[2]{H(z, z^*)}{q}{z_i}
	\end{equation}
	
	Here we must consider the real embedding of $z$ to properly analyse the emergent dynamics. Using $z_i = x_i + iy_i$ we obtain
	\begin{equation}
	\left(\pdv[2]{H(z, z^*)}{z_i}{z_j} + \pdv[2]{H(z, z^*)}{z_i}{z_j^*}\right) \pdv{x_j}{q} + \left(\pdv[2]{H(z, z^*)}{z_i}{z_j} - \pdv[2]{H(z, z^*)}{z_i}{z_j^*}\right) i \pdv{y_j}{q} = -\pdv[2]{H(z, z^*)}{q}{z_i}
	\end{equation}
	which can be reduced into a system of real equations
	\begin{align}
	\Re\left(\pdv[2]{H(z, z^*)}{z_i}{z_j} + \pdv[2]{H(z, z^*)}{z_i}{z_j^*}\right) \pdv{x_j}{q} - \Im\left(\pdv[2]{H(z, z^*)}{z_i}{z_j} - \pdv[2]{H(z, z^*)}{z_i}{z_j^*}\right)\pdv{y_j}{q} &= -\Re\pdv[2]{H(z, z^*)}{q}{z_i}\\
	\Im\left(\pdv[2]{H(z, z^*)}{z_i}{z_j} + \pdv[2]{H(z, z^*)}{z_i}{z_j^*}\right) \pdv{x_j}{q} + \Re\left(\pdv[2]{H(z, z^*)}{z_i}{z_j} - \pdv[2]{H(z, z^*)}{z_i}{z_j^*}\right)\pdv{y_j}{q} &= -\Im\pdv[2]{H(z, z^*)}{q}{z_i}
	\end{align}
	
	We define a general state vector
	
	\begin{equation}
	\zeta = \mqty(x_1 \\ \vdots \\ x_N \\ y_i \\ \vdots \\ y_N)
	\end{equation}
	
	Then the equation of motion permits this useful form
	
	\begin{equation}
	\Omega \dv{\zeta}{q} = R\qq{where}\Omega = \mqty(\Omega^{RR} & \Omega^{RI} \\ \Omega^{IR} & \Omega^{II})\qq{and}R = \mqty(R^R \\ R^I)
	\end{equation}
	where
	\begin{align}
	\Omega^{RR}_{ij} &= \Re\left(\pdv[2]{H(z, z^*)}{z_i}{z_j} + \pdv[2]{H(z, z^*)}{z_i}{z_j^*}\right)\\
	\Omega^{RI}_{ij} &= - \Im\left(\pdv[2]{H(z, z^*)}{z_i}{z_j} - \pdv[2]{H(z, z^*)}{z_i}{z_j^*}\right)\\
	\Omega^{IR}_{ij} &= \Im\left(\pdv[2]{H(z, z^*)}{z_i}{z_j} + \pdv[2]{H(z, z^*)}{z_i}{z_j^*}\right)\\
	\Omega^{II}_{ij} &= \Re\left(\pdv[2]{H(z, z^*)}{z_i}{z_j} - \pdv[2]{H(z, z^*)}{z_i}{z_j^*}\right)
	\end{align}
	and
	\begin{align}
	R^R_i &= -\Re\pdv[2]{H(z, z^*)}{q}{z_i}\\
	R^I_i &= -\Im\pdv[2]{H(z, z^*)}{q}{z_i}
	\end{align}
	This uniquely specifies the characteristic trajectory of $\ket{z}$ under adiabatic evolution (assuming $\Omega$ is non-singular).
	
	\textit{Note.} Anticipating the next subsection, it is also helpful to solve for $\pdv{z_j}, \pdv{z_j^*}$ directly without the real embedding. For this, we consider the conjugate stationary condition, which is equivalent to the complex conjugate of the expanded equation of motion:
	\begin{equation}
	\pdv[2]{H(z, z^*)}{z_i^*}{z_j^*} \pdv{z_j^*}{q} + \pdv[2]{H(z, z^*)}{z_i^*}{z_j} \pdv{z_j}{q} = -\pdv[2]{H(z, z^*)}{q}{z_i^*}
	\end{equation}
	where we use the fact that $q$ is real and $H(z, z^*)$ is real-valued when $z^*$ is actually taken to be equal to the complex conjugate of $z$ (i.e. when the equation represents the expectation energy value of $\ket{z}$). The resulting system of equations is
	\begin{equation}
	\mqty(\Omega^{AA} & \Omega^{AB} \\ \Omega^{BA} & \Omega^{BB}) \mqty(\pdv{z^*}{q} \\ \pdv{z}{q}) = \mqty(R^A \\ R^B)
	\end{equation}
	where
	\begin{align}
	\Omega^{AA}_{ij} &= \pdv[2]{H(z, z^*)}{z_i^*}{z_j^*} = \alpha_{ij}\\
	\Omega^{AB}_{ij} &= \pdv[2]{H(z, z^*)}{z_i^*}{z_j} = \beta_{ij}\\
	\Omega^{BA}_{ij} &= \pdv[2]{H(z, z^*)}{z_i}{z_j^*} = \beta^*_{ij}\\
	\Omega^{BB}_{ij} &= \pdv[2]{H(z, z^*)}{z_i}{z_j} = \alpha^*_{ij}\\
	R^A_i &= -\pdv[2]{H(z, z^*)}{q}{z_i^*} \qq{and} R^B_i = -\pdv[2]{H(z, z^*)}{q}{z_i} = R^{A*}
	\end{align}
	Then by standard identities:
	\begin{equation}
	\mqty(\Omega^{AA} & \Omega^{AB} \\ \Omega^{BA} & \Omega^{BB})\inv = \mqty((\alpha - \beta \alpha^{-1*} \beta^*)\inv & 0 \\ 0 & (\alpha^* - \beta^* \alpha\inv \beta)\inv) \mqty(I & -\beta \alpha^{-1*} \\ -\beta^* \alpha\inv & I)
	\end{equation}
	which yields the useful form
	\begin{align} \label{eq: characteristic trajectories no embedding}
	\pdv{z}{q} &= (\alpha^* - \beta^* \alpha\inv \beta)\inv (-\beta^* \alpha\inv R^A + R^B)\\
	\pdv{z^*}{q} &=  (\alpha - \beta \alpha^{-1*} \beta^*)\inv (R^A - \beta \alpha^{-1*} R^B)
	\end{align}
	Note that this approach requires $\Omega^{AA}$ and $\Omega^{BB}$ to be invertible and, for the purpose of a stable numerical calculation, well-conditioned.
	
	
	\subsubsection{Variational equations of motion for a composite state}
	
	In our ansatz, a general state will be approximated by a finite decomposition
	\begin{equation}
	\ket{\Psi} = \sum_1^N c_a \unnormket{z_a}
	\end{equation}
	Constraining the $q$-dependence of $\unnormket{z_a}$ to its characteristic trajectory, the dynamical properties of $\ket{\Psi}$ are only the decomposition coefficients $c_a$. Therefore, the full equations of motion are
	\begin{align}
	\pdv{c_a^*} \dv{F}{q} &= 0\\
	\pdv{c_a} \dv{F}{q} &= 0
	\end{align}
	where we substitute $\dot{z}_{a,j}^*, \dot{z}_{b,j}$ from Eqns.~\eqref{eq: characteristic trajectories no embedding}. We also make the substitution
	\begin{align}
	\pdv[2]{H(z_a, z_a^*)}{z_{a,i}^*}{z_{a,j}^*} = \alpha_{a,ij} = W_{aa}(\hat{T}_i\hc \hat{T}_j\hc \hat{H}) / S_{aa}\\
	\pdv[2]{H(z_a, z_a^*)}{z_{a,i}^*}{z_{a,j}} = \beta_{a,ij} = W_{aa}(\hat{T}_i\hc \hat{H} \hat{T}_j) / S_{aa}
	\end{align}
	where $S$ is the overlap matrix:
	\begin{equation}
	S_{ab} = \unnormbraket{z_a}{z_b} = W_{ab}(\hat{I})
	\end{equation}
	The equations of motion can then be expanded as
	\begin{align}
	\mqty(\Gamma^{AA} & \Gamma^{AB} \\ \Gamma^{BA} & \Gamma^{BB}) \dv{q}\mqty(c^* \\ c) = \mqty(Q^A \\ Q^B)
	\end{align}
	where
	\begin{align}
	\Gamma^{AA}_{mn} &= -2W_{ab}(\hat{H})W_{cd}(\hat{H}) \delta_c^m \delta_a^n c_b c_d\\
	\Gamma^{AB}_{mn} &= \delta_a^m \delta_b^n W_{ab}(\hat{H}^2) - 2W_{ab}(\hat{H})W_{cd}(\hat{H})(\delta_a^m  c_c^* + \delta_c^m c_a^*) \delta_b^n c_d\\
	\Gamma^{BA}_{mn} &= \delta_b^m \delta_a^n W_{ab}(\hat{H}^2) - 2W_{ab}(\hat{H}) W_{cd}(\hat{H}) (\delta_b^m c_d + \delta_d^m c_b) \delta_a^n c_c^*\\
	\Gamma^{BB}_{mn} &= -2W_{ab}(\hat{H})W_{cd}(\hat{H}) \delta_d^m \delta_b^n c_a^* c_c^*\\
	Q^A_m &= -2 \left( \delta^m_a c_b c_c^* c_d + \delta^m_c c_a^* c_b c_d \right) W_{cd}(\hat{H}) \left[ W_{ab}\left(\pdv{q}\hat{H}\right) + \left( W_{ab}(\hat{T}_j\hc \hat{H}) \dot{z}_{a,j}^* + W_{ab}(\hat{H} \hat{T}_j) \dot{z}_{b,j} \right) \right] \nonumber\\
	&\quad + \delta^m_a c_b \left[  W_{ab}\left(\pdv{\hat{H}^2}{q}\right) + W_{ab}(\hat{T}_j\hc \hat{H}^2) \dot{z}_{a,j}^* + W_{ab}(\hat{H}^2 \hat{T}_j) \dot{z}_{b,j}\right]\\
	Q^B_m &= -2 \left( \delta^m_b c_a^* c_c^* c_d + \delta^m_d c_a^* c_b c_c^* \right) W_{cd}(\hat{H}) \left[ W_{ab}\left(\pdv{q}\hat{H}\right) + \left( W_{ab}(\hat{T}_j\hc \hat{H}) \dot{z}_{a,j}^* + W_{ab}(\hat{H} \hat{T}_j) \dot{z}_{b,j} \right) \right] \nonumber\\
	&\quad +\delta^m_b c_a^* \left[  W_{ab}\left(\pdv{\hat{H}^2}{q}\right) + W_{ab}(\hat{T}_j\hc \hat{H}^2) \dot{z}_{a,j}^* + W_{ab}(\hat{H}^2 \hat{T}_j) \dot{z}_{b,j}\right]
	\end{align}
	where we have the expected relations
	\begin{equation}
	{\Gamma^{BB}_{mn}}^* = \Gamma^{AA}_{mn} \qquad {\Gamma^{BA}_{mn}}^* = \Gamma^{AB}_{mn} \qquad {Q^B_m}^* = Q^A_m
	\end{equation}
	Then the system of equations may be expanded as
	\begin{equation}
	\Gamma^{AA} \dot{c}^* + \Gamma^{AB} \dot{c} = Q^A \qquad {\Gamma^{AB}}^* \dot{c}^* + {\Gamma^{AA}}^* \dot{c} = {Q^A}^*
	\end{equation}
	from which we have the substitution
	\begin{equation}
	\dot{c}^* = {\Gamma^{AB}}^{*-1} ({Q^A}^* - {\Gamma^{AA}}^* \dot{c})
	\end{equation}
	yielding
	\begin{equation}
	(\Gamma^{AB} - \Gamma^{AA} {\Gamma^{AB}}^{*-1} {\Gamma^{AA}}^*) \dot{c} = Q^A - \Gamma^{AA} {\Gamma^{AB}}^{*-1} {Q^A}^*
	\end{equation}
	Based on Sec. TODO
	
	\section{Why the long face}
	In a perfect world, every property of coherent states could be treated using the ansatz of differential geometry. For example, adiabatic transfer characteristric trajectories could be described using some special curvature $\mathcal{C}$ and scalar $\mathcal{F}$ for which
	\begin{equation}
	\dv{z_i}{q} = \mathcal{C}_{ij} \pdv{F}{z_j^*}
	\end{equation}
	where $\mathcal{C}, \mathcal{T}$ are functions of $H(z^*, z)$ and $\dv{H}{q}$.
	
	...but this isn't a perfect world...
	
	By eliminating in different order we get
	
	\begin{equation}
	\pdv{z}{q} = G_{ij}^{-1} D_j
	\end{equation}
	where
	\begin{align}
	G &= \beta - \alpha \beta^* \alpha^{*}\\
	D &= R^A - \alpha \beta^{-1*} R^B
	\end{align}
	Note that $G$ isa Hermitian bilinear form, but not a Kahler metric since $\pdv{z_k}G_{ij}$ is not invariant under index rotation (check)
	\begin{align}
	\pdv{z_k}G_{ij} &= 
	\end{align}
	
	
	
	
\end{document}